\begin{center}
    \textbf{\large Вариант 1}
\end{center}

\samerowans{\begin{problem}{1}
Вычислите: $5709 - 5632 : (258 - 13 \cdot 19) + 103$
\end{problem}}

\vspace{5pt}

\samerowans{\begin{problem}{2}
Решите уравнение: $(294 : a - 14 \cdot 6) : 3 = 21$
\end{problem}}

\vspace{5pt}

\samerowans{\begin{problem}{3}
Вычислите: $9\text{ т } 54\text{ кг } - 76\text{ ц } + 127\text{ кг}$. Ответ запишите в т, ц, кг.
\end{problem}}

\samerowans{\begin{problem}{4}
Сложили пять одинаковых чисел, к результату прибавили 8, затем полученную сумму разделили на 9 и получили 12. Какие числа складывали?
\end{problem}}

\vspace{5pt}

\samerowans{\begin{problem}{5}
Оля два часа делала домашнее задание. Три четверти часа она делала математику, полчаса русский. Все оставшееся время она учила английский. Сколько минут Оля учила английский?
\end{problem}}

\vspace{5pt}

\begin{problem}{6} Из двух пунктов А и В, расстояние между которыми 50 км, одновременно навстречу друг другу выехали два велосипедиста. Скорость велосипедиста, выехавшего из п А, на $5\text{ км/ч}$ больше скорости второго. За два часа им удалось встретиться и разъехаться на расстояние 40 км. На каком расстоянии от п А будет второй велосипедист через 3 часа 15 минут после своего выезда из В?
\end{problem}

\vspace{5pt}

\noindent\grid{36}{19}

\vspace{5pt}

\samerowans{\begin{problem}{7}
Первый станок может сделать 300 деталей за 6 часов, второй справится с этой задачей в 2 раза быстрее. Сколько часов совместной работы потребуется двум станкам, чтобы сделать 750 деталей?
\end{problem}}

\newpage

\begin{problem}{8} Из прямоугольника с периметром 16 дм вырезали квадрат со стороной 5 см. Найди площадь оставшейся фигуры, если длина прямоугольника 500 мм.
\end{problem}

\smallskip

\noindent\grid{36}{15}

\vspace{5pt}

\samerowans{\begin{problem}{9}
Таня заплатила за покупку 1 м 50 см шелковой ленты 600 рублей. Надя купила ленты на 80 см больше, чем Таня. Сколько рублей заплатила Надя?
\end{problem}}

\vspace{5pt}

\samerowans{\begin{problem}{10}
Коля и Петя сыграли партию в шахматы, которая закончилась победой одного из них. «Я победил», — сказал Коля. «Я проиграл», — сказал Петя. Кто из них выиграл, если известно, что хотя бы кто-то из них сказал неправду?
\end{problem}}

\vspace{5pt}

\begin{problem}{11} Через два года мальчик будет вдвое старше, чем он был два года назад. А девочка будет через три года втрое старше, чем три года назад. Кто старше: мальчик или девочка?
\end{problem}

\smallskip

\noindent\grid{36}{20}

\newpage


\begin{center}
\textsc{\large Дополнительный лист}
\end{center}
\noindent\grid{34}{47}

\newpage

\null
